\subsubsection{Euler tour / Hierholzer (camino y circuito)}

\paragraph{Idea intuitiva.}
Encuentra un \textbf{camino o circuito euleriano} (que usa cada arista exactamente una vez). \textbf{Condiciones}:
\begin{itemize}\setlength\itemsep{0pt}
	\item \textbf{No dirigido}: todos los grados pares (circuito) o exactamente 2 impares (camino, empieza en uno, termina en el otro). Ademas, conexo sobre el subgrafo con aristas.
	\item \textbf{Dirigido}: $\text{indeg}(v) = \text{outdeg}(v)$ para todos (circuito) o exactamente un nodo con $\text{outdeg} = \text{indeg} + 1$ y otro con $\text{indeg} = \text{outdeg} + 1$ (camino).
\end{itemize}

\paragraph{Hierholzer}: elegir un ciclo desde el inicio, mientras haya ciclos ``laterales'' desde un nodo del camino, mergearlos. Implementacion con pila: avanzar por aristas, al no poder mas aniadir a la respuesta y backtrack. La demostracion de $O(V + E)$: cada arista se procesa una vez (entra y sale del stack una vez).

\paragraph{Propiedades clave (invariantes).}
\begin{itemize}\setlength\itemsep{0pt}
	\item Multigrafos permitidos; las aristas se marcan como usadas.
	\item No dirigido: cada arista aparece \textbf{dos veces} en \texttt{adj} (ida y vuelta).
	\item La construccion del path con ids consistentes es importante.
	\item El resultado es unico salvo orden de las aristas en multigrafos.
\end{itemize}

\paragraph{Codigo clave.}
Hierholzer con pila:
\begin{verbatim}
vector<int> path;
stack<int> st;
st.push(start);
while (!st.empty()) {
    int v = st.top();
    while (!adj[v].empty() && used[adj[v].back()]) adj[v].pop_back();
    if (adj[v].empty()) {
        path.push_back(v);
        st.pop();
    } else {
        int e = adj[v].back();
        used[e] = true;
        int u = edges[e].to(v);
        adj[v].pop_back();
        st.push(u);
    }
}
reverse(path.begin(), path.end());
\end{verbatim}

\paragraph{Complejidad.}
\begin{itemize}\setlength\itemsep{0pt}
	\item $O(V + E)$.
	\item Memoria $O(V + E)$.
\end{itemize}

\paragraph{Donde tocar para modificar.}
\begin{itemize}\setlength\itemsep{0pt}
	\item \textbf{Reconstruir con multigrafo}: aniadir ids consistentes antes del algoritmo.
	\item \textbf{Aplicar a ``de Bruijn'' o problemas de strings}: modelar como multigrafo y buscar Euler.
	\item \textbf{Dirigido vs no dirigido}: el snippet tiene ambas versiones.
	\item \textbf{Letter addition}: modelar palabras como aristas y buscar el superstring.
\end{itemize}

\paragraph{Trampas y errores frecuentes.}
\begin{itemize}\setlength\itemsep{0pt}
	\item Si el grafo no es conexo (sobre las aristas), no existe euleriano. Verificar primero.
	\item Para multigrafos, los IDs de aristas son cruciales para distinguirlas.
	\item Si el problema es ``encuentra todos los euler tours'', el algoritmo es exponencial.
\end{itemize}

\paragraph{Codigo.}
Ver \texttt{cpp.tex}, seccion \textit{Euler tour / Hierholzer (camino y circuito)}.

\vspace{0.5em}
